\documentclass{beamer}

\usepackage[utf8]{inputenc}
\usepackage[english]{babel}
\usepackage{textpos}
\usepackage{graphicx}
\usepackage{textcomp}
\usepackage{animate}
\usepackage{comment}
%\usepackage{multicol}
%\setlength{\columnsep}{5cm}
\usepackage{amsmath}
\usepackage{amsfonts}
\usepackage{verbatim}
\usepackage{varwidth}
\usepackage{caption}
\usepackage{dcolumn}
\usepackage{amssymb}
\usepackage{algorithm}
\usepackage{algorithmic}
%\usepackage{physics}
\usepackage{hyperref}
\usepackage{scalerel,stackengine}
%\usepackage{fancyvrb}


\newcommand{\code}[1]{\texttt{#1}}
\newcommand{\shp}{$>>>$ }
\newcommand{\tab}{\hspace*{22pt}}

\usetheme{Madrid}

\title{Python in One Line}
\author{EEC 3150}
\date{}

\institute{Trevecca Nazarene University}
 


%%%%%%%%% BEGIN DOCUMENT %%%%%%%%%
%%%%%%%%% BEGIN DOCUMENT %%%%%%%%%
%%%%%%%%% BEGIN DOCUMENT %%%%%%%%%

\begin{document}

% get title page 
\frame{\titlepage}
 
%%%%%%%%%%
%%%%%%%%%%
\begin{frame}	
	\frametitle{Outline}
	\tableofcontents	
\end{frame}


%%%%%%%%%%
%%%%%%%%%%
\section{Inline If Statements}

\begin{frame}	
	\center{\Huge{Inline If Statements}}
\end{frame}

\begin{frame}	
	\frametitle{Inline If Statements}
	
	\begin{block}{Inline if statements:}	
		You can write simple if statements in one line in the following way: \\[2pt]
		\code{ans = x if \emph{condition} else y}
		
		\vspace{10pt}
		Examples:
		\begin{itemize}
			\item \code{x = 10\\
				y = 0 if x\%2==0 else 1
		}
			\item \code{n = 3 \\
				x = 'a' if n==0 else 'b' if n==1 else 'c'
		}
		\end{itemize}
		\vspace{10pt}
		You can't do inline \code{elif} but you can achieve the same result by nesting inline ifs
	\end{block}
	
\end{frame}


%%%%%%%%%%
%%%%%%%%%%
\section{List Comprehension}

\begin{frame}	
	\center{\Huge{List Comprehension}}
\end{frame}

\begin{frame}	
	\frametitle{List Comprehension}
	
	\begin{block}{List Comprehension:}	
		You can create lists and dictionaries in one line using comprehension: \\[2pt]
		\code{myList = [ \emph{expr} for i in \emph{iterable} ]} \\[2pt]
		\code{myDict = [ \emph{key}:\emph{value} for i in \emph{iterable} ]} \\[2pt]
		
		\vspace{10pt}
		Examples:
		\begin{footnotesize}
		\begin{itemize}
			\item \code{x = [ i for i in range(10) ]} (integers 0-9)
			\item \code{x = [ i for i in range(10) if i\%2==0 ]} (even integers 0-9)
			\item \code{x = [ y[i] for i in range(len(y)) if y[i]>thresh ]} (filter another list)
			\item \code{x = [ [0 for j in range(5)] for i in range(5) ]} (5x5 of 0's)
		\end{itemize}
		\end{footnotesize}
	\end{block}
	
\end{frame}




\end{document}




